\documentclass{IEEEcsmag}

% Add this package for superscript citations (optional, but requested by template text)
\usepackage[super]{cite}

\usepackage[colorlinks,urlcolor=blue,linkcolor=blue,citecolor=blue]{hyperref}
\expandafter\def\expandafter\UrlBreaks\expandafter{\UrlBreaks\do\/\do\*\do\-\do\~\do\'\do\"\do\-}
\usepackage{upmath,color}


\jvol{XX}
\jnum{XX}
\paper{8}
\jmonth{Month}
\jname{Publication Name}
\jtitle{Publication Title}
\pubyear{2021}

\newtheorem{theorem}{Theorem}
\newtheorem{lemma}{Lemma}


\setcounter{secnumdepth}{0}

\begin{document}

\sptitle{Article Type: Description  (see below for more detail)}

\title{The Reproducibility Crisis of LLM Benchmarks in Software Engineering: An Industrial Perspective}

\author{Gustavo A. Oliva}
\affil{Queen's University, Kingston, ON, K7L 3N6, Canada}

\author{Dayi Lin}
\affil{Queen's University, Kingston, ON, K7L 3N6, Canada}

\author{Ahmed E. Hassan}
\affil{Queen's University, Kingston, ON, K7L 3N6, Canada}

\markboth{THEME/FEATURE/DEPARTMENT}{THEME/FEATURE/DEPARTMENT}

\begin{abstract}\looseness-1The integration of Large Language Models (LLMs) into Software Engineering (SE) has spurred the creation of benchmarks like SWE-Bench and LCB to measure their capabilities. However, current practices for reporting results are failing us, leading to a reproducibility crisis. The execution of these benchmarks depends on a vast number of parameters—from the LLM's training cutoff and inference settings to the specific versions of evaluation scaffolds—many of which go unreported. While some parameters have a greater impact than others, any one of them can alter the outcome, making it impossible for industrial teams to establish reliable baselines for their own models. This paper argues that the SE community must move beyond isolated, single-point benchmark publications. Drawing inspiration from mature practices in hardware benchmarking, we propose a paradigm shift towards a crowdsourced and continuously updated evaluation model. In this model, benchmark results are submitted by various parties, complete with detailed configuration reports, to establish a public, transparent distribution of performance data (e.g., median, average). This approach would not only expose the impact of different configurations but also create the robust, trustworthy baselines that industry desperately needs. We contend that such a shift is essential for fostering genuine progress and enabling fair comparison in the evaluation of LLMs for software engineering.
\end{abstract}

\maketitle

The rapid advancement of Large Language Models (LLMs) has created a paradigm shift in software engineering (SE), with models now capable of generating code, fixing bugs, and even designing entire software components. To measure their progress, the community has developed benchmarks like SWE-Bench, LiveCodeBench, and the Aider Polyglot benchmark. However, the current practices for reporting results are failing, leading to a crisis in reproducibility, transparency, and trustworthiness.

From an industrial perspective, this is a critical obstacle. As a company training our own LLMs, we rely on public benchmarks to establish a performance baseline. Furthermore, executing these benchmarks is a time-consuming and computationally expensive endeavor. Integrating multiple benchmarks, each with its own unique setup and non-standard parameters, is akin to a complex legacy system integration project. This lack of standardization leads to wasted resources and makes it impossible to reliably compare results. For instance, a public, crowdsourced baseline would allow us to see how far a new model's official scores deviate from the median, providing a much-needed reality check.

In this paper, we distill our team's practical experience into a set of five concrete recommendations for the community. We begin with a motivating example, using our own data to show how minor parameter variations cause significant swings in benchmark scores. We then present our core recommendations: (1) Mandate cryptographic hashes of model weights for verifiable integrity. (2) Require a comprehensive, machine-readable configuration file for all runs. (3) Enforce the reporting of a "default settings" baseline alongside any optimized results. (4) Champion a shift to a crowdsourced, statistical benchmarking model. (5) Issue a call for standardized benchmark engineering, treating benchmarks as robust, containerized tools with standard interfaces and resilient execution. Our goal is to move the community towards a more robust, reliable, and scientific foundation for evaluating LLMs in software engineering.

% TODO: Focus on industrial context of building LLMs, but also mention this is all valid for academics/researchers creating agentic scaffolds for solving SWE tasks (e.g., alternatives to OpenHands, etc).

\section{MOTIVATING EXAMPLE}
\label{sec:motivation}

% The idea here is to show one execution of SWEBench with OpenHands. We will use default OpenHands parameters. Then, we'll show a second execution where only one change in the OpenHands parameters drastically improves the results. We will also show that OpenHands documentation actually disencourages activating said parameter.

\section{RECOMMENDATIONS}
\label{sec:recommendations}

In this section, we present three concrete recommendations derived from our industrial experience. These recommendations are designed to be practical and actionable, addressing the key challenges of reproducibility, transparency, and trustworthiness in LLM benchmarking for software engineering.

\subsection{Comprehensive and Standardized Configuration Reporting}
\label{sec:rec1}
A fundamental barrier to reproducibility is the ad-hoc and incomplete manner in which benchmark configurations are currently reported. To address this, we propose a standardized, multi-part reporting structure that covers all layers of the evaluation stack. This structure should be captured in a \texttt{benchmark\_spec.yaml} file to ensure clarity and consistency.

\subsubsection{Model Specification}
Trust in benchmark results begins with knowing precisely which model was evaluated. Currently, submissions often only specify a model's name (e.g., "Llama-3-70B"). This is insufficient, as fine-tuned variants can have vastly different capabilities. We propose that all submissions must specify the exact base model and include a cryptographic hash (e.g., SHA-256) of the model weights for any open-source model. This provides an unambiguous, verifiable fingerprint of the exact artifact used.

% TODO: Add Qwen3 Coder example where metadata doesn't align with published model.

\subsubsection{Inference Engine Parameters}
The software used to serve the model can significantly impact performance. Parameters related to the deployment and optimization of the inference engine, such as the configuration of vLLM or TensorRT-LLM, are often overlooked. This includes settings like the level of quantization, data-parallel size, and other engine-specific tunings. These parameters must be explicitly reported to understand the full context of the execution environment.

% TODO: Add our VLLM parameter set

\subsubsection{Inference Parameters}
The parameters used at the time of generation directly control the model's output behavior. While some benchmarks report basic settings like temperature, a complete picture is necessary. We recommend the mandatory reporting of all inference parameters, including but not limited to temperature, top-p, top-k, repetition penalty, and maximum token count. 

% TODO: List the 30B models that we investigated and the key set of inference parameters that we identified.

\subsubsection{Scaffold Parameters}
The "scaffold" or "harness" that orchestrates the interaction between the model and the benchmark tasks is a major source of variability. Frameworks like OpenHands for SWE-Bench have their own configuration parameters that can dramatically influence outcomes. These settings, which control aspects like the agent's strategy or the number of interaction turns, must be treated as first-class citizens in the configuration report.

% TODO: Add our reference parameters for Aider, LiveCodeBench, and SWE-Bench.

\subsubsection{Execution Environment}
When a benchmark is not executed within a provided container, the host environment becomes a significant source of variability. We recommend that all non-containerized benchmark runs must report a complete snapshot of the execution environment. This includes the operating system, Python version, and the specific version of critical hardware drivers like CUDA. Furthermore, a full list of pinned library dependencies, such as the output of \texttt{pip freeze}, must be included to prevent discrepancies arising from package updates.

\subsection{Shift to Crowdsourced, Statistical Benchmarking}
\label{sec:rec2}

The current model of relying on a few single-point results published in papers or by model vendors is fragile and prone to cherry-picking. We advocate for a paradigm shift towards a crowdsourced, statistical model, inspired by mature fields like hardware benchmarking. We envision a central, public repository where the community can submit benchmark results, each accompanied by its detailed configuration file (as per Recommendation 2). This would create a living, transparent distribution of performance data for any given model, allowing for the calculation of robust statistical measures like median and average performance. Such a system would serve as a powerful "reality check" against official claims and provide a much more reliable signal of a model's true capabilities.

Model providers, in their quest to showcase the best possible performance, often publish benchmark scores obtained from highly-tuned configurations. While these "hero run" results are valuable for demonstrating a model's peak potential, they can be misleading if presented in isolation. We recommend that any benchmark result published by a model provider must be accompanied by a result from a run using the benchmark's official, out-of-the-box default scaffold parameters. This mandatory baseline provides a stable, common reference point, allowing the community to fairly assess the true "alpha" generated by custom tuning versus the model's core capabilities under a standardized setting.

\subsection{Standardized Benchmark Engineering}
\label{sec:rec3}
From an industrial perspective, integrating and running multiple benchmarks is a significant software engineering challenge. Each benchmark is often a standalone, idiosyncratic system, making scaled evaluation a complex and costly integration project. We issue a call for the community to treat benchmarks as robust, interoperable software tools. This includes standardizing their interfaces, such as supporting a common set of inference parameters. It also means engineering for resilience and efficiency, for example, by immediately serializing the result of each task to prevent catastrophic data loss from a single failure. Finally, benchmarks should be distributed as containerized applications (e.g., using Docker) to ensure they can be run in a safe, isolated, and easily manageable environment.

% \section{Related Work}
% \label{sec:related-work}

% Related work critiquing benchmarks practices in ML/AI
% Discuss EvalHub type of solutions

\section{Conclusion}
\label{sec:conclusion}

% TODO once paper is ready

[The conclusion]

Fake citation: \cite{goth2024polyglot}

% % TODO: Praise the SE community for creating so many benchmarks and the effort that has gone into them. Say we live in this transformative era of agentic AI (Cite Ahmed's paper) and we as a community must make the best of out of these amazing tools.

\section{SUPPLEMENTARY MATERIAL}

See the Appendix, which is available in the IEEE Computer Society Digital Library at \url{http://doi.ieeecomputersociety.org/10.1109/MCSE.2019.2947017}.

For detailed algorithms, please refer to the supplementary materials, which can be found in the Computer Society Digital Library at \url{http://doi.ieeecomputersociety.org/10.1109/MCSE.2019.2947017}.\vspace*{-8pt}

\section{ACKNOWLEDGMENTS}
The author(s) would like to thank A, B, and C. This work was supported by XYZ under Grant \#\#\#.

\def\refname{REFERENCES}

\bibliographystyle{IEEEtran}
\bibliography{references}


\begin{IEEEbiography}{Example} {\,} is a researcher at the  ABC Corporation, B\"oblingen, Germany.  Her current research interests include a, b, and c. Author received her Ph.D. degree  in physics from University. She is a Fellow of the IEEE Computer Society. Contact her at sbauthor@abc.com.
\end{IEEEbiography}

\begin{IEEEbiography}{Gustavo A. Oliva}{\,}is ...
\end{IEEEbiography}

\begin{IEEEbiography}{Dayi Lin}{\,} is 
\end{IEEEbiography}

\begin{IEEEbiography}{Ahmed E. Hassan} {\,} is 
\end{IEEEbiography}

\end{document}

